\begin{table}[htbp]\centering\footnotesize
\caption{El perímetro pensionario por institución, y el Ramo 19 bruto y neto, línea P (mdp)}\label{cua:pensinst}
\begin{tabular}{>{\raggedright\arraybackslash}p{4.6cm}rrrr}
\toprule
Institución & 2026 & 2027 & Dif. nominal & Var. real (\%) \\
\midrule
IMSS & 1.010.525,2 & 1.136.228,0 & 125.702,8 & +8,1 \\
ISSSTE & 395.156,0 & 409.867,0 & 14.711,0 & $-$0,3 \\
Pemex & 92.333,4 & 94.472,9 & 2.139,5 & $-$1,6 \\
CFE & 67.841,7 & 66.154,8 & $-$1.686,9 & $-$6,2 \\
Resto del Gobierno Federal, por diferencia & 138.302,9 & 134.023,4 & $-$4.279,5 & $-$6,8 \\
\midrule \textbf{Perímetro del Anexo 3} & \textbf{1.704.159,2} & \textbf{1.840.746,1} & \textbf{136.586,9} & \textbf{+3,9} \\
\midrule \multicolumn{5}{l}{\emph{Ramo 19: no se suma a lo anterior}} \\
Ramo 19 bruto & 1.541.518,7 & 1.661.007,9 & 119.489,2 & +3,6 \\
Transferido a IMSS e ISSSTE (UR GYR y GYN) & 1.401.952,8 & 1.526.114,6 & 124.161,8 & +4,7 \\
Ramo 19 neto de esas transferencias & 139.565,9 & 134.893,2 & $-$4.672,7 & $-$7,1 \\
\bottomrule
\end{tabular}
\\[2pt]\begin{minipage}{0.92\textwidth}\scriptsize Fuente: analíticos del proyecto 2026 y 2027, tipo de gasto 4 (pensiones y jubilaciones), \texttt{datos/ruta\_a/pensiones\_tg4\_institucion.csv}; perímetro por diferencia de los dos renglones del Anexo 3 de los decretos. Deflactor del PIB 1,040. \textbf{El Ramo 19 bruto cuenta dos veces lo que transfiere al IMSS y al ISSSTE}, que reaparece en el presupuesto de esas entidades; el neto es la cifra que publica CIEP.\end{minipage}
\end{table}
